% GENERATED FILE. Edit canonical lessons, then run npm run generate:books.
% canonical-source-hash: fnv1a64:cdf7f7f79b7fbf7a
% canonical-lessons: HI-C488-jac, HI-C488-nuskha, HI-C488-nars, HI-C488-dikhvana, HI-C488-karana

\chapter{At the Doctor's}
\label{ch:hi-at-the-doctor-s}
\hlchaptermodality{488}

\begin{chapteropening}
\textbf{By the end of this chapter:} \emph{I can talk about a visit to the doctor.}

Complete the last lesson of chapter 488: I can talk about a visit to the doctor.
\end{chapteropening}

\section[\texorpdfstring{\hi{जाँच} (jā̃c)}{जाँच (jā̃c)}]{\hi{जाँच} (jā̃c) — a check-up; a test}

\label{lesson:HI-C488-jac}



\begin{quote}
Before the new one: say the Hindi for smallpox, then the Hindi for to fester.
\end{quote}

\subsection*{You'll want to know: \hi{जाँच}}
\textbf{\hi{जाँच}} — \emph{jā̃c} — \textquotedblleft{}a check-up; a test\textquotedblright{}.

\begin{quote}
\hi{डाक्टर} \hi{ने} \hi{मीरा} \hi{की} \hi{आँखों} \hi{की} \hi{जाँच} \hi{की।}
\end{quote}

\begin{quote}
\emph{ḍākṭar ne Mīrā kī ā̃khoṁ kī jā̃c kī.}
\end{quote}

\begin{quote}
The doctor examined Mira's eyes.
\end{quote}

Feminine.

\begin{cousinweb}
\textbf{\hi{जाँच}} (\emph{jā̃c}) is inherited: CDIAL lists Hindi \emph{jā̃c}, \textquotedblleft{}examination\textquotedblright{}, under Sanskrit \emph{yācñā}, \textquotedblleft{}request, asking\textquotedblright{} (CDIAL 10448): asking questions of a thing. \emph{jā̃c karnā} is to examine.
\end{cousinweb}

\begin{grammarlens}[title={What you've built}]
One more word for describing what you lived and felt.
\end{grammarlens}

\subsection*{Your turn}
\begin{itemize}
\raggedright
  \item \emph{Say it:} \emph{jā̃c}
  \item \emph{Say it:} \emph{jā̃c}, once more
  \item \emph{Say it:} \emph{ḍākṭar ne Mīrā kī ā̃khoṁ kī jā̃c kī.} — The doctor examined Mira's eyes.
  \item \emph{Recall:} say the Hindi for smallpox, then the Hindi for to fester, then say \emph{jā̃c} again
\end{itemize}

\begin{tcolorbox}[breakable,colback=teal!4,colframe=teal!35!black,title={Before you move on}]
What does \hi{जाँच} mean? (\textquotedblleft{}A check-up; a test\textquotedblright{}.) Say it once more.
\end{tcolorbox}
\section[\texorpdfstring{\hi{नुस्ख़ा} (nuskhā)}{नुस्ख़ा (nuskhā)}]{\hi{नुस्ख़ा} (nuskhā) — a prescription}

\label{lesson:HI-C488-nuskha}



\begin{quote}
Before the new one: say the Hindi for to fester, then the Hindi for a check-up.
\end{quote}

\subsection*{You'll want to know: \hi{नुस्ख़ा}}
\textbf{\hi{नुस्ख़ा}} — \emph{nuskhā} — \textquotedblleft{}a prescription\textquotedblright{}.

\begin{quote}
\hi{डाक्टर} \hi{ने} \hi{नया} \hi{नुस्ख़ा} \hi{लिखा।}
\end{quote}

\begin{quote}
\emph{ḍākṭar ne nayā nuskhā likhā.}
\end{quote}

\begin{quote}
The doctor wrote out a new prescription.
\end{quote}

Masculine.

\begin{cousinweb}
\textbf{\hi{नुस्ख़ा}} (\emph{nuskhā}) is usually traced to Arabic \emph{nuskha}, \textquotedblleft{}a copy, a written text\textquotedblright{}, which reached Hindi through Persian; a recipe for medicine is a written text. It is also a household remedy.
\end{cousinweb}

\begin{grammarlens}[title={What you've built}]
One more word for describing what you lived and felt.
\end{grammarlens}

\subsection*{Your turn}
\begin{itemize}
\raggedright
  \item \emph{Say it:} \emph{nuskhā}
  \item \emph{Say it:} \emph{nuskhā}, once more
  \item \emph{Say it:} \emph{ḍākṭar ne nayā nuskhā likhā.} — The doctor wrote out a new prescription.
  \item \emph{Recall:} say the Hindi for to fester, then the Hindi for a check-up, then say \emph{nuskhā} again
\end{itemize}

\begin{tcolorbox}[breakable,colback=teal!4,colframe=teal!35!black,title={Before you move on}]
What does \hi{नुस्ख़ा} mean? (\textquotedblleft{}A prescription\textquotedblright{}.) Say it once more.
\end{tcolorbox}
\section[\texorpdfstring{\hi{नर्स} (nars)}{नर्स (nars)}]{\hi{नर्स} (nars) — a nurse}

\label{lesson:HI-C488-nars}



\begin{quote}
Before the new one: say the Hindi for a check-up, then the Hindi for a prescription.
\end{quote}

\subsection*{You'll want to know: \hi{नर्स}}
\textbf{\hi{नर्स}} — \emph{nars} — \textquotedblleft{}a nurse\textquotedblright{}.

\begin{quote}
\hi{नर्स} \hi{ने} \hi{मरीज़} \hi{को} \hi{दवा} \hi{दी।}
\end{quote}

\begin{quote}
\emph{nars ne marīz ko davā dī.}
\end{quote}

\begin{quote}
The nurse gave the patient his medicine.
\end{quote}

Masculine or feminine, by the person meant.

\begin{cousinweb}
\textbf{\hi{नर्स}} (\emph{nars}) is borrowed from English \emph{nurse}.
\end{cousinweb}

\begin{grammarlens}[title={What you've built}]
One more word for describing what you lived and felt.
\end{grammarlens}

\subsection*{Your turn}
\begin{itemize}
\raggedright
  \item \emph{Say it:} \emph{nars}
  \item \emph{Say it:} \emph{nars}, once more
  \item \emph{Say it:} \emph{nars ne marīz ko davā dī.} — The nurse gave the patient his medicine.
  \item \emph{Recall:} say the Hindi for a check-up, then the Hindi for a prescription, then say \emph{nars} again
\end{itemize}

\begin{tcolorbox}[breakable,colback=teal!4,colframe=teal!35!black,title={Before you move on}]
What does \hi{नर्स} mean? (\textquotedblleft{}A nurse\textquotedblright{}.) Say it once more.
\end{tcolorbox}
\section[\texorpdfstring{\hi{दिखवाना} (dikhvānā)}{दिखवाना (dikhvānā)}]{\hi{दिखवाना} (dikhvānā) — to have (something) looked at}

\label{lesson:HI-C488-dikhvana}



\begin{quote}
Before the new one: say the Hindi for a prescription, then the Hindi for a nurse.
\end{quote}

\subsection*{You'll want to know: \hi{दिखवाना}}
\textbf{\hi{दिखवाना}} — \emph{dikhvānā} — \textquotedblleft{}to have (something) looked at\textquotedblright{}.

\begin{quote}
\hi{अरुण} \hi{ने} \hi{अपनी} \hi{रिपोर्ट} \hi{डाक्टर} \hi{को} \hi{दिखवाई।}
\end{quote}

\begin{quote}
\emph{Aruṇ ne apnī riporṭ ḍākṭar ko dikhvāī.}
\end{quote}

\begin{quote}
Arun had the doctor look at his report.
\end{quote}

Past, with \hi{ने} (he or she …): \hi{उसने} \hi{आँख} \hi{दिखवाई}, \emph{usne ā̃kh dikhvāī}.

\begin{cousinweb}
\textbf{\hi{दिखवाना}} (\emph{dikhvānā}) is the second causative of \emph{dekhnā}, \textquotedblleft{}to see\textquotedblright{}, which CDIAL lists under \emph{*dṛkṣati}, a form it reconstructs (CDIAL 6507): \emph{dikhānā} is \textquotedblleft{}to show\textquotedblright{}, and \emph{dikhvānā} is \textquotedblleft{}to have something seen to\textquotedblright{}, usually by a doctor.
\end{cousinweb}

\begin{grammarlens}[title={What you've built}]
One more word for describing what you lived and felt.
\end{grammarlens}

\subsection*{Your turn}
\begin{itemize}
\raggedright
  \item \emph{Say it:} \emph{dikhvānā}
  \item \emph{Say it:} \emph{dikhvānā}, once more
  \item \emph{Say it:} \emph{Aruṇ ne apnī riporṭ ḍākṭar ko dikhvāī.} — Arun had the doctor look at his report.
  \item \emph{Recall:} say the Hindi for a prescription, then the Hindi for a nurse, then say \emph{dikhvānā} again
\end{itemize}

\begin{tcolorbox}[breakable,colback=teal!4,colframe=teal!35!black,title={Before you move on}]
What does \hi{दिखवाना} mean? (\textquotedblleft{}To have (something) looked at\textquotedblright{}.) Say it once more.
\end{tcolorbox}
\section[\texorpdfstring{\hi{कराना} (karānā)}{कराना (karānā)}]{\hi{कराना} (karānā) — to have something done}

\label{lesson:HI-C488-karana}



\begin{quote}
Before the new one: say the Hindi for a nurse, then the Hindi for to have (something) looked at.
\end{quote}

\subsection*{You'll want to know: \hi{कराना}}
\textbf{\hi{कराना}} — \emph{karānā} — \textquotedblleft{}to have something done\textquotedblright{}.

\begin{quote}
\hi{अरुण} \hi{ने} \hi{अस्पताल} \hi{में} \hi{ख़ून} \hi{की} \hi{जाँच} \hi{कराई।}
\end{quote}

\begin{quote}
\emph{Aruṇ ne aspatāl meṁ khūn kī jā̃c karāī.}
\end{quote}

\begin{quote}
Arun had a blood test done at the hospital.
\end{quote}

Past, with \hi{ने} (he or she …): \hi{उसने} \hi{जाँच} \hi{कराई}, \emph{usne jā̃c karāī}.

\begin{cousinweb}
\textbf{\hi{कराना}} (\emph{karānā}) is the causative partner of \emph{karnā}, \textquotedblleft{}to do\textquotedblright{}: CDIAL lists Hindi \emph{karānā}, \textquotedblleft{}to cause to be made\textquotedblright{}, under Sanskrit \emph{kar\'{\={o}}ti}, \textquotedblleft{}does\textquotedblright{} (CDIAL 2814). Someone else does the work for you.
\end{cousinweb}

\begin{grammarlens}[title={What you've built}]
One more word for describing what you lived and felt.
\end{grammarlens}

\subsection*{Your turn}
\begin{itemize}
\raggedright
  \item \emph{Say it:} \emph{karānā}
  \item \emph{Say it:} \emph{karānā}, once more
  \item \emph{Say it:} \emph{Aruṇ ne aspatāl meṁ khūn kī jā̃c karāī.} — Arun had a blood test done at the hospital.
  \item \emph{Recall:} say the Hindi for a nurse, then the Hindi for to have (something) looked at, then say \emph{karānā} again
\end{itemize}

\begin{tcolorbox}[breakable,colback=teal!4,colframe=teal!35!black,title={Before you move on}]
What does \hi{कराना} mean? (\textquotedblleft{}To have something done\textquotedblright{}.) Say it once more.
\end{tcolorbox}
